%=======================================================================
% TCC v4 - PARTE 2: MATERIAIS E MÉTODOS + SISTEMA PROPOSTO
% Gabriel Figueiredo Carolino - Sistema de Irrigação Automatizada IoT
%=======================================================================

%=======================================================================
% 3 MATERIAIS E MÉTODOS
%=======================================================================
\chapter{Materiais e Métodos}

\section{Metodologia Científica}

Este trabalho adota a abordagem metodológica de \textbf{Design Science Research} (DSR), caracterizada pelo desenvolvimento de artefatos tecnológicos para resolver problemas práticos identificados, seguindo rigor científico na construção, demonstração e avaliação da solução proposta \cite{hevner2004}. A metodologia DSR é particularmente adequada para pesquisas em sistemas de informação e engenharia que visam criar soluções inovadoras para problemas reais, combinando criatividade no design com validação empírica rigorosa.

A pesquisa foi estruturada seguindo as seis etapas fundamentais da DSR: (1) \textbf{Identificação do problema} - caracterização das dificuldades no cuidado de plantas domésticas e limitações das soluções existentes; (2) \textbf{Definição de objetivos} - estabelecimento de requisitos claros para uma solução que equilibre funcionalidade técnica com simplicidade de uso; (3) \textbf{Design e desenvolvimento} - criação da arquitetura IoT integrada com hardware, software e interfaces; (4) \textbf{Demonstração} - implementação de protótipo funcional operando em cenários realísticos; (5) \textbf{Avaliação} - validação através de testes técnicos e de experiência do usuário; (6) \textbf{Comunicação} - documentação detalhada de resultados e contribuições para a comunidade científica.

A validação experimental segue princípios de \textbf{pesquisa experimental aplicada}, utilizando metodologia quantitativa para métricas técnicas de desempenho e abordagem quali-quantitativa para avaliação de experiência do usuário. Esta combinação metodológica permite avaliar tanto a eficácia técnica quanto a aceitação prática da solução desenvolvida.

A justificativa para esta abordagem metodológica reside na natureza eminentemente prática do problema abordado - o cuidado automatizado de plantas domésticas - que demanda soluções que funcionem efetivamente no mundo real, não apenas em ambientes controlados de laboratório. A DSR proporciona o rigor científico necessário enquanto mantém foco na aplicabilidade prática dos resultados.

% [INSERIR IMAGEM 3: Esquema da metodologia DSR aplicada ao projeto]
% Inserir aqui um diagrama mostrando as 6 etapas da Design Science Research
% aplicadas especificamente ao desenvolvimento do sistema de irrigação

\section{Arquitetura do Sistema e Abordagem Experimental}

O sistema foi projetado seguindo uma arquitetura em três camadas que equilibra processamento local com capacidades de nuvem, garantindo operação confiável mesmo durante intermitências de conectividade. Esta arquitetura modular facilita tanto a replicação quanto a extensão do sistema para diferentes contextos de aplicação.

% [INSERIR IMAGEM 4: Arquitetura de camadas do sistema completo]
% Inserir aqui um diagrama detalhado mostrando:
% - Camada Física (ESP32, sensores, atuadores)
% - Camada de Comunicação (Wi-Fi, Firebase)
% - Camada de Interface (Web responsiva, controles adaptativos)

A \textbf{camada física} compreende os componentes de hardware responsáveis pela interação direta com o ambiente: microcontrolador ESP32 como unidade central de processamento, sensores de umidade do solo (capacitivo), temperatura e umidade do ar (DHT22), luminosidade (LDR), e sistema de acionamento de bomba d'água através de circuito amplificador com transistor PNP.

A \textbf{camada de comunicação} gerencia a troca de dados entre o sistema local e serviços em nuvem, implementando protocolos de reconexão automática, buffer local para operação offline e sincronização inteligente de dados via Firebase Realtime Database.

A \textbf{camada de interface} oferece acesso ao sistema através de interface web responsiva que se adapta automaticamente a diferentes dispositivos, implementando dois modos operacionais distintos (simples e avançado) para atender usuários com diferentes níveis de experiência técnica.

% [INSERIR IMAGEM 5: Fluxograma do experimento de validação]
% Inserir aqui um fluxograma mostrando:
% - Configuração inicial (3 vasos diferentes)
% - Processo de calibração por vaso
% - Período de monitoramento (14 dias por vaso)
% - Coleta de dados técnicos e de experiência do usuário
% - Análise e validação dos resultados

\section{Especificação de Hardware e Materiais}

\subsection{Componentes Principais}

O microcontrolador \textbf{ESP32 DevKit v1} foi selecionado como plataforma principal devido à combinação ótima de capacidades de processamento, conectividade integrada e custo acessível. O modelo utilizado oferece 30 pinos GPIO, processamento dual-core de 240 MHz, 520 KB de SRAM, conectividade Wi-Fi 802.11 b/g/n e Bluetooth, proporcionando recursos suficientes para operação simultânea de múltiplas tarefas críticas.

Para sensoriamento de umidade do solo, foi utilizado \textbf{sensor capacitivo v1.2} resistente à corrosão, conectado ao pino analógico GPIO34 (ADC1_CH6) que oferece resolução de 12 bits (4096 níveis) para leituras precisas. A escolha do sensor capacitivo sobre alternativas resistivas foi motivada pela maior durabilidade, estabilidade temporal e ausência de corrosão eletrolítica em contato prolongado com substrato úmido.

O monitoramento climático é realizado através de sensor \textbf{DHT22} (AM2302), oferecendo medição simultânea de temperatura (-40°C a +80°C com precisão ±0.5°C) e umidade relativa (0-100% com precisão ±2% RH). O protocolo de comunicação digital one-wire simplifica a integração e reduz susceptibilidade a interferências eletromagnéticas.

Para medição de luminosidade, foi implementado sensor \textbf{LDR (Light Dependent Resistor)} GL5528 em configuração de divisor de tensão, oferecendo resposta logarítmica adequada para a ampla faixa de intensidades luminosas típicas em ambientes internos (10 lux a 1000 lux).

\subsection{Sistema de Acionamento e Alimentação}

O controle da bomba de irrigação utiliza arquitetura de amplificação de corrente baseada em transistor \textbf{TIP127 (PNP)} para acionamento de módulo relé 5V/10A. Esta configuração foi selecionada para garantir isolamento galvânico adequado entre circuitos de controle e potência, além de capacidade de corrente suficiente para bombas submersíveis de pequeno porte.

% [INSERIR IMAGEM 6: Esquema elétrico completo do circuito]
% Inserir aqui o diagrama de circuito mostrando:
% - ESP32 com identificação de todos os pinos utilizados
% - Conexões dos sensores (umidade, DHT22, LDR)
% - Circuito do transistor TIP127 + relé + bomba
% - Fonte de alimentação 5V
% - Conexões de terra (GND comum)

O circuito implementa lógica invertida apropriada para transistores PNP, onde um sinal LOW no pino GPIO4 do ESP32 ativa o transistor, energizando o relé e acionando a bomba. Um resistor de 1kΩ limita a corrente de base do transistor (I_base ≈ 4.3mA), garantindo operação segura dentro das especificações de ambos os componentes (I_base_max = 5mA para ESP32, I_base_típica = 5mA para TIP127).

A bomba d'água selecionada (modelo submersível 3-6V, 120L/h) opera com baixo consumo de corrente (≤ 300mA), adequada para aplicações domésticas com vasos de 1 a 10 litros. A escolha por bomba de baixa tensão contribui para segurança operacional e simplifica o sistema de alimentação, eliminando necessidade de conversores adicionais.

\section{Metodologia de Calibração Experimental}

\subsection{Protocolo de Calibração por Volume de Vaso}

A calibração dos sensores de umidade constitui aspecto metodologicamente crítico desta pesquisa, especialmente considerando as diferenças significativas entre vasos de diferentes volumes. O protocolo foi desenvolvido para ser executado de forma sistemática e reproduzível, estabelecendo base científica sólida para operação confiável do sistema.

O processo de calibração seguiu metodologia experimental controlada com as seguintes etapas:

\textbf{Etapa 1 - Preparação do substrato}: Utilização de solo universal comercial padronizado (composição: 60% turfa, 20% vermiculita, 15% casca de pinus, 5% aditivos), peneirado em malha 5mm para garantir homogeneidade granulométrica. Volume específico por vaso: 0,8L (pequeno), 2,5L (médio), 6,0L (grande).

\textbf{Etapa 2 - Determinação do ponto seco}: Secagem controlada em ambiente ventilado (umidade relativa 40-50%, temperatura 23±2°C) por 48 horas, com medições a cada 4 horas para garantir estabilização completa. Critério de estabilização: variação ≤ 0,1% entre medições consecutivas. O valor médio das últimas 10 medições define o parâmetro \textit{rawDry}.

\textbf{Etapa 3 - Determinação do ponto saturado}: Saturação controlada do substrato através de irrigação gradual (50ml a cada 5 minutos) até atingir capacidade de campo, definida como cessação de drenagem livre após 30 minutos. Após período de equilíbrio de 2 horas, medições estabilizadas definem o parâmetro \textit{rawWet}.

\textbf{Etapa 4 - Validação da linearidade}: Medições em pontos intermediários (25%, 50%, 75% de saturação) para verificar linearidade da resposta sensor. Cálculo do coeficiente de correlação linear (R²) para validação da qualidade da calibração.

\subsection{Controle de Variáveis Externas}

Para garantir validade científica dos resultados experimentais, foram implementadas medidas rigorosas de controle de variáveis externas que poderiam influenciar as medições:

\textbf{Controle térmico}: Todos os experimentos foram conduzidos em ambiente com temperatura controlada (22±2°C) através de ar condicionado, com monitoramento contínuo via sensor DHT22. Variações térmicas podem afetar tanto a condutância dos sensores capacitivos quanto a taxa de evapotranspiração das plantas.

\textbf{Controle de luminosidade}: Ambiente experimental mantido com iluminação artificial estável (lâmpadas LED 6500K, 40W, posicionadas a 50cm dos vasos) durante período diurno (6h às 18h), eliminando variações de luz natural que poderiam criar gradientes de evaporação diferentes entre vasos.

\textbf{Controle de umidade relativa}: Manutenção da umidade do ar entre 50-60% através de umidificador ultrassônico controlado por timer, evitando que variações atmosféricas influenciem diferencialmente a evaporação conforme tamanho dos vasos.

\textbf{Controle de correntes de ar}: Posicionamento dos vasos em área sem correntes de ar diretas, protegida de ventilação de ar condicionado ou janelas, garantindo que a evaporação ocorra uniformemente e seja dependente apenas das características intrínsecas de cada configuração de vaso.

\textbf{Controle de substrato}: Utilização do mesmo lote de solo comercial para todos os vasos, com homogeneização prévia através de mistura manual e peneiramento, eliminando variações na composição que poderiam afetar propriedades dielétricas e retenção hídrica.

\section{Ferramentas de Desenvolvimento e Análise}

\subsection{Ambiente de Desenvolvimento}

O desenvolvimento do firmware para ESP32 foi realizado utilizando \textbf{Arduino IDE 2.2.1} com bibliotecas específicas: WiFi.h (conectividade), Firebase\_ESP\_Client v4.4.12 (comunicação com banco de dados), e bibliotecas auxiliares para tokenização e manipulação RTDB. A escolha do ambiente Arduino foi motivada pela maturidade do ecossistema, disponibilidade de bibliotecas e facilidade de replicação por outros pesquisadores.

A interface web foi desenvolvida em \textbf{HTML5, CSS3 e JavaScript ES6+}, utilizando Firebase SDK v10.7.1 para comunicação em tempo real com o banco de dados. Para visualização de dados, foi integrada a biblioteca \textbf{Chart.js v4.4.4}, oferecendo gráficos interativos responsivos que se adaptam automaticamente a diferentes tamanhos de tela.

\subsection{Ferramentas de Coleta e Análise de Dados}

A coleta de dados experimentais foi estruturada utilizando múltiplas ferramentas para garantir redundância e facilitar análise posterior:

\textbf{Coleta primária}: Firebase Realtime Database com estrutura hierárquica organizada por dispositivo, tipo de dados (snapshot, telemetria, eventos, configuração) e timestamp. Frequência de coleta: 30 segundos para dados ambientais, eventos síncronos para ações de irrigação.

\textbf{Backup local}: Buffer circular no ESP32 (últimas 100 medições) para garantir continuidade da coleta durante interrupções de conectividade, com sincronização automática quando conexão é restabelecida.

\textbf{Exportação para análise}: Scripts Python 3.9 para extração de dados do Firebase em formato CSV, facilitando análise estatística posterior em planilhas eletrônicas e ferramentas especializadas.

\textbf{Análise estatística}: Cálculo de médias, desvios-padrão, coeficientes de variação e testes de normalidade utilizando \textbf{Microsoft Excel 2021} e \textbf{Python pandas/numpy} para análises mais complexas. Métricas de desempenho incluem precisão de medição (comparação com referências), responsividade temporal (latência entre detecção e ação) e confiabilidade operacional (uptime e taxa de falhas).

\section{Protocolo de Teste e Validação}

\subsection{Cenário Experimental Controlado}

O protocolo experimental foi desenvolvido para validar a eficácia do sistema em condições que simulam fielmente o uso doméstico real, considerando variações comuns de tamanhos de vaso e espécies de plantas encontradas em residências urbanas.

\textbf{Configuração de vasos}: Três vasos representativos foram selecionados: (1) pequeno - 1,2L, diâmetro 12cm, altura 10cm, para plantas como suculentas e ervas; (2) médio - 3,5L, diâmetro 18cm, altura 15cm, para plantas ornamentais de porte médio; (3) grande - 7,8L, diâmetro 25cm, altura 20cm, para plantas maiores ou pequenas árvores domésticas.

\textbf{Espécies vegetais}: Seleção de plantas com necessidades hídricas distintas para validar adaptabilidade do sistema: manjericão (Ocimum basilicum) no vaso pequeno - alta demanda hídrica; samambaia de Boston (Nephrolepis exaltata) no vaso médio - demanda moderada; espada-de-são-jorge (Sansevieria trifasciata) no vaso grande - baixa demanda hídrica.

\textbf{Duração e periodicidade}: Teste de longo prazo com 14 dias para cada configuração de vaso (total: 42 dias), permitindo validação de estabilidade temporal e adaptação a variações sazonais naturais. Coleta de dados contínua 24h/dia com armazenamento automático.

% [INSERIR IMAGEM 7: Setup experimental mostrando os três vasos instrumentados]
% Inserir aqui uma foto do ambiente experimental completo:
% - Três vasos de diferentes tamanhos com plantas
% - ESP32 e sensores instalados
% - Sistema de irrigação montado
% - Ambiente controlado com iluminação artificial

\subsection{Métricas de Avaliação Técnica e de Experiência do Usuário}

A avaliação do sistema foi estruturada considerando tanto aspectos técnicos quantificáveis quanto fatores subjetivos relacionados à experiência do usuário, proporcionando validação abrangente da solução desenvolvida.

\textbf{Métricas técnicas de desempenho}:
\begin{itemize}
\item \textbf{Precisão de medição}: Erro médio absoluto e desvio-padrão das leituras de umidade comparadas com medições gravimétricas de referência (balança analítica ±0.1g)
\item \textbf{Responsividade temporal}: Latência entre detecção de condição de irrigação e acionamento efetivo da bomba, incluindo tempos de comunicação Firebase
\item \textbf{Confiabilidade operacional}: Percentual de disponibilidade do sistema, tempo médio entre falhas (MTBF) e tempo médio de recuperação (MTTR)
\item \textbf{Eficiência energética}: Consumo de energia total do sistema durante operação normal e modos de baixo consumo
\end{itemize}

\textbf{Métricas de experiência do usuário}:
\begin{itemize}
\item \textbf{Facilidade de configuração inicial}: Tempo necessário para usuários não técnicos completarem instalação e calibração básica
\item \textbf{Transparência operacional}: Capacidade dos usuários compreenderem e confiarem nas decisões automatizadas do sistema
\item \textbf{Satisfação geral}: Avaliação através de questionários estruturados baseados no Technology Acceptance Model (TAM)
\item \textbf{Intenção de uso continuado}: Indicadores de aderência e recomendação da solução para outros usuários
\end{itemize}

A validação através do \textbf{Technology Acceptance Model (TAM)} seguiu metodologia estabelecida na literatura de sistemas de informação, avaliando percepção de utilidade, facilidade de uso, intenção comportamental e uso efetivo através de questionários padronizados aplicados após diferentes períodos de exposição ao sistema (1 semana, 2 semanas, 1 mês).

%=======================================================================
% 4 SISTEMA PROPOSTO
%=======================================================================
\chapter{Sistema Proposto}

\section{Visão Geral da Arquitetura}

O sistema proposto representa uma solução integrada de automação residencial que combina simplicidade de uso com robustez técnica, seguindo princípios de design centrado no usuário sem comprometer capacidades técnicas essenciais. A arquitetura foi projetada considerando modularidade, escalabilidade e facilidade de manutenção, garantindo que o sistema possa ser facilmente replicado, adaptado e estendido para diferentes contextos domésticos.

A filosofia central do projeto baseia-se no conceito de "automação transparente", onde a tecnologia opera de forma quase invisível para o usuário, fornecendo cuidados contínuos e inteligentes às plantas sem exigir intervenção ou conhecimento técnico constante. Esta abordagem reconhece que o valor principal da automação residencial reside na eliminação de preocupações e tarefas repetitivas, não na exibição de complexidade tecnológica.

% [INSERIR IMAGEM 8: Arquitetura geral do sistema em diagrama de blocos]
% Inserir aqui um diagrama de alto nível mostrando:
% - Três subsistemas principais e suas interações
% - Fluxos de dados bidirecionais
% - Pontos de decisão automática
% - Interfaces de comunicação (Wi-Fi, Firebase, Web)

O sistema integra três subsistemas principais que operam de forma coordenada e síncrona: o \textbf{subsistema de monitoramento ambiental}, responsável pela coleta contínua de dados sobre condições do solo e ambiente circundante; o \textbf{subsistema de controle inteligente}, que processa informações coletadas e toma decisões automatizadas sobre irrigação baseadas em algoritmos adaptativos; e o \textbf{subsistema de interface e comunicação}, que proporciona acesso remoto, transparência operacional e controle manual quando necessário.

\section{Subsistema de Monitoramento Ambiental}

O subsistema de monitoramento ambiental constitui a base sensorial do sistema, implementando estratégia de coleta de dados multi-paramétrica que vai além da simples medição de umidade do solo. Esta abordagem holística reconhece que as necessidades hídricas das plantas são influenciadas por múltiplas variáveis ambientais que interagem de forma complexa.

\subsection{Sensoriamento de Umidade do Solo}

O sensor de umidade capacitivo representa o componente mais crítico do subsistema, posicionado estrategicamente no substrato para capturar variações representativas da disponibilidade hídrica na zona radicular. O posicionamento considera a distribuição típica do sistema radicular da espécie cultivada e evita áreas próximas às bordas do vaso onde a evaporação periférica pode causar leituras não representativas do estado geral do substrato.

A tecnologia capacitiva foi selecionada por sua estabilidade temporal superior, ausência de corrosão eletrolítica e resposta linear em ampla faixa de umidade. O sensor opera medindo variações na constante dielétrica do substrato, que correlaciona diretamente com o conteúdo de água devido às propriedades elétricas distintivas da molécula H₂O.

\subsection{Monitoramento Climático Integrado}

O monitoramento climático através do sensor DHT22 fornece dados essenciais sobre temperatura e umidade do ar, variáveis que influenciam significativamente a taxa de transpiração das plantas e, consequentemente, suas necessidades hídricas. O sistema implementa algoritmos que correlacionam dados climáticos com medições de umidade do solo para otimizar decisões de irrigação.

Temperaturas elevadas e baixa umidade do ar aumentam a demanda transpirativa, exigindo ajustes nos limiares de irrigação. Conversely, condições de alta umidade relativa e temperaturas amenas reduzem esta demanda, permitindo que o sistema postergue irrigações sem comprometer o bem-estar das plantas.

\subsection{Sensor de Luminosidade e Fotoperíodo}

O sensor de luminosidade LDR complementa o monitoramento ao fornecer informações sobre intensidade e duração da radiação disponível para fotossíntese. Plantas expostas a maior luminosidade apresentam metabolismo mais ativo, maior taxa fotossintética e, consequentemente, maiores necessidades hídricas.

O sistema utiliza dados de luminosidade para implementar ajustes circadianos nos algoritmos de controle, reconhecendo que plantas têm necessidades hídricas diferentes durante períodos de atividade fotossintética diurna versus respiração noturna. Esta funcionalidade avançada diferencia a solução de sistemas mais simples que operam com limiares fixos independentes do ciclo natural das plantas.

\section{Subsistema de Controle Inteligente}

O subsistema de controle inteligente implementa a lógica central que transforma dados sensoriais em ações de irrigação contextualmente apropriadas. A arquitetura de controle baseia-se em uma máquina de estados finitos que opera autonomamente no microcontrolador ESP32, garantindo operação confiável mesmo durante interrupções de conectividade ou falhas de comunicação.

\subsection{Máquina de Estados e Lógica de Decisão}

A máquina de estados implementa três estados principais com transições bem definidas:

\textbf{IDLE (Monitoramento Passivo)}: Estado normal de operação onde o sistema monitora continuamente condições ambientais e avalia necessidade de intervenção. A transição para irrigação ocorre quando umidade do solo cruza abaixo do limiar inferior configurado, considerando também dados contextuais de temperatura, umidade do ar e luminosidade.

\textbf{IRRIGATING (Irrigação Ativa)}: Durante este estado, o sistema ativa a bomba d'água e monitora simultaneamente umidade do solo e tempo decorrido. A irrigação continua até que umidade atinja limiar superior configurado ou até que tempo máximo de segurança seja atingido, implementando proteção contra irrigação excessiva.

\textbf{LOCKOUT (Período de Estabilização)}: Estado de espera obrigatório após cada ciclo de irrigação, permitindo distribuição uniforme da água no substrato e estabilização das leituras sensoriais. Este período previne oscilações indesejadas e garante que decisões subsequentes sejam baseadas em dados estabilizados.

% [INSERIR IMAGEM 9: Diagrama da máquina de estados com transições]
% Inserir aqui um fluxograma detalhado mostrando:
% - Os três estados principais (IDLE, IRRIGATING, LOCKOUT)
% - Condições específicas para cada transição
% - Parâmetros monitorados em cada estado
% - Proteções de segurança implementadas

\subsection{Algoritmos Adaptativos e Proteções de Segurança}

Além da lógica de estados básica, o sistema implementa múltiplas camadas de proteção e adaptação:

\textbf{Proteções temporais}: Tempo máximo absoluto de irrigação (configurável, padrão 60s), intervalos mínimos entre ciclos (configurável, padrão 15 minutos), timeout para detecção de sensores desconectados (30 segundos).

\textbf{Proteções de valores}: Detecção de anomalias nos sensores (valores fora de faixas esperadas), validação de leituras através de médias móveis, rejeição de spikes causados por interferência eletromagnética.

\textbf{Adaptação contextual}: Ajuste automático de limiares baseado em histórico de comportamento da planta, consideração de fatores climáticos na tomada de decisão, aprendizado das características específicas de cada configuração de vaso/substrato.

\textbf{Modo de segurança}: Capacidade de operação em modo degradado quando dados críticos não estão disponíveis, utilizando parâmetros conservadores para garantir sobrevivência das plantas.

\section{Subsistema de Interface e Comunicação}

O subsistema de interface representa a ponte entre automação técnica e experiência do usuário, implementando filosofia de design que prioriza simplicidade e intuitividade sem sacrificar funcionalidade avançada para usuários experientes.

\subsection{Interface Web Responsiva e Adaptativa}

A interface web foi desenvolvida utilizando tecnologias modernas (HTML5, CSS3, JavaScript ES6+) que garantem compatibilidade com ampla gama de dispositivos e navegadores. O design responsivo adapta-se automaticamente a smartphones, tablets e computadores desktop, oferecendo experiência consistente independente da plataforma de acesso.

A arquitetura de interface implementa conceito inovador de \textbf{modos adaptativos}:

\textbf{Modo Simples}: Utiliza controles visuais intuitivos como sliders, botões grandes e indicadores gráficos. Usuários configuram quando suas plantas devem ser regadas usando linguagem natural ("solo seco", "solo úmido") em vez de percentuais técnicos. Funcionalidades avançadas ficam ocultas, reduzindo complexidade cognitiva.

\textbf{Modo Avançado}: Oferece acesso completo a todos os parâmetros técnicos, incluindo calibração manual de sensores, configuração detalhada de algoritmos de controle, análise de dados históricos e exportação de logs para análise externa.

Esta dualidade garante que o sistema atenda tanto usuários iniciantes (que valorizam simplicidade) quanto experientes (que desejam controle granular) sem comprometer a experiência de nenhum grupo.

% [INSERIR IMAGEM 10: Interface web mostrando modo simples vs avançado]
% Inserir aqui uma comparação lado a lado:
% - Modo simples: sliders, linguagem natural, controles grandes
% - Modo avançado: campos numéricos, gráficos técnicos, configurações detalhadas

\subsection{Comunicação em Tempo Real e Persistência de Dados}

A comunicação via Firebase Realtime Database implementa arquitetura push-based que garante sincronização instantânea entre dispositivo IoT e interfaces de usuário. Esta sincronização bidirecional permite tanto monitoramento remoto quanto controle em tempo real.

\textbf{Estrutura de dados hierárquica}:
\begin{itemize}
\item \texttt{/devices/\{deviceId\}/snapshot}: Estado atual (umidade, status bomba, timestamp)
\item \texttt{/devices/\{deviceId\}/config}: Configuração ativa (limiares, calibração, parâmetros)
\item \texttt{/devices/\{deviceId\}/telemetry}: Histórico temporal para análise de tendências
\item \texttt{/devices/\{deviceId\}/commands}: Comandos remotos (irrigação manual, reconfiguração)
\item \texttt{/devices/\{deviceId\}/events}: Log de eventos para auditoria e debugging
\end{itemize}

\textbf{Funcionalidades de histórico e análise}: Gráficos interativos mostram variações de umidade, frequência de irrigação e correlações com condições ambientais ao longo do tempo. Usuários podem identificar padrões, otimizar configurações e detectar problemas precocemente através de análise visual intuitiva.

\section{Adaptação para Diferentes Tamanhos de Vasos}

Uma das principais inovações do sistema é sua capacidade de adaptação automática para vasos de diferentes volumes, reconhecendo que esta variação é ubíqua em ambientes domésticos e afeta fundamentalmente as características de distribuição de umidade e resposta à irrigação.

\subsection{Metodologia de Calibração Específica}

\textbf{Vasos pequenos (≤ 2L)}: Configurações otimizadas para resposta rápida e distribuição uniforme de umidade. Algoritmos implementam ciclos de irrigação mais curtos e frequentes, considerando menor inércia térmica e maior susceptibilidade a variações ambientais.

\textbf{Vasos médios (2-5L)}: Parâmetros intermediários que equilibram responsividade com estabilidade. Ciclos de irrigação de duração moderada com intervalos apropriados para permitir distribuição adequada sem saturação excessiva.

\textbf{Vasos grandes (> 5L)}: Reconhecimento de distribuição heterogênea de umidade e resposta mais lenta. Algoritmos ajustados para ciclos mais longos com intervalos estendidos, garantindo penetração adequada da água em todo o volume de substrato.

% [INSERIR IMAGEM 11: Comparação dos três tamanhos de vasos utilizados]
% Inserir aqui uma foto mostrando:
% - Três vasos lado a lado (pequeno, médio, grande)
% - Identificação clara das capacidades (1L, 3L, 7L)
% - Posicionamento correto dos sensores em cada vaso
% - Diferentes espécies de plantas utilizadas nos testes

\subsection{Algoritmos de Compensação Volumétrica}

O sistema implementa algoritmos matemáticos que compensam automaticamente diferenças volumétricas:

\textbf{Compensação temporal}: Ajuste da duração de irrigação baseado no volume do vaso (t\_irrigação = k × √volume, onde k é constante específica da bomba utilizada).

\textbf{Compensação de limiares}: Ajuste automático dos percentuais de umidade considerando diferentes características de distribuição (vasos grandes requerem limiares ligeiramente mais altos devido à heterogeneidade espacial).

\textbf{Compensação de intervalos}: Cálculo automático de períodos de lockout baseado na cinética de distribuição de umidade (vasos maiores requerem períodos mais longos para estabilização).

Esta abordagem sistêmica garante precisão operacional independente do tamanho do vaso, eliminando necessidade de usuários compreenderem ou ajustarem manualmente estes parâmetros complexos.

\section{Características de Segurança e Confiabilidade}

O sistema incorpora múltiplas camadas de proteção operando em níveis complementares (hardware, firmware, aplicação) para garantir operação segura em ambiente doméstico desacompanhado.

\subsection{Proteções de Hardware}

\textbf{Isolamento galvânico}: Módulo relé com optoacoplador fornece isolamento elétrico completo entre circuitos de controle (5V) e carga (bomba), eliminando risco de danos por sobrecorrente ou curto-circuito.

\textbf{Limitação de corrente}: Resistores de precisão limitam corrente de base do transistor amplificador, garantindo operação dentro de especificações tanto do ESP32 quanto do TIP127.

\textbf{Proteção contra inversão}: Diodo de proteção em paralelo com a bobina do relé previne picos de tensão indutiva que poderiam danificar componentes semicondutores.

\subsection{Proteções de Firmware}

\textbf{Watchdog timer}: Reinicialização automática do sistema em caso de travamento de software, garantindo recuperação automática de falhas não previstas.

\textbf{Tratamento de exceções}: Captura e tratamento sistemático de erros de comunicação, leitura de sensores e operações matemáticas, evitando comportamentos indefinidos.

\textbf{Buffer circular}: Armazenamento local das últimas configurações válidas, permitindo operação autônoma mesmo durante falhas prolongadas de conectividade.

\textbf{Logs detalhados}: Registro de todas as operações críticas com timestamp, facilitando diagnóstico de problemas e análise de desempenho histórico.

A integração dessas múltiplas camadas de proteção resulta em sistema robusto capaz de operar de forma confiável e segura em ambiente doméstico típico, mesmo com usuários não técnicos e sem supervisão constante.
